% !TEX TS-program = pdflatexmk
\documentclass[letter,12pt]{article} %[twocolumn]
% page formatting %%%%%%%%%%%%%%%%%%%%%
%\usepackage{geometry} %defaults are roughly [top=33mm, bottom=48mm, left=33mm, right=33mm]
%\usepackage[top=20mm, bottom=30mm, left=15mm, right=15mm]{geometry}
\usepackage[top=30mm, bottom=40mm, left=25mm, right=25mm]{geometry}
%\usepackage[top=30mm, bottom=45mm, left=30mm, right=30mm]{geometry}
%\usepackage[top=26mm, bottom=40mm, left=26mm, right=26mm]{geometry}

\usepackage{/Users/wdlinch3/Dropbox/Rashoumon/LaTeX/macros/WDL3}

\usepackage{/Users/wdlinch3/Dropbox/Rashoumon/LaTeX/macros/SupergeometryDefs}
\usepackage{/Users/wdlinch3/Dropbox/Rashoumon/LaTeX/macros/SupergeometryDefs2}


%%%%%%%%%%%%%%%%%%%%%%%%%%%%%%%%%%%%%%%%%%%%%%%%%%%%%%%%%%%%%%%%%
%%%%%%%%%%%%%%%%%%%%%%%%%%%%%%%%%%%%%%%%%%%%%%%%%%%%%%%%%%%%%%%%%
%%%%%%%%%%%%%%%%%%%%%%%%%%%%%%%%%%%%%%%%%%%%%%%%%%%%%%%%%%%%%%%%%
%Definitions for this paper 

\usepackage{xcolor}
\usepackage{pagecolor}

\def\gauge#1{ {\color{green} { {#1}} } }
\def\physical#1{ {\color{green} {\bm {#1}} } }
\def\auxiliary#1{ {\color{yellow} {\bm {#1}} } }
\def\compensating#1{ {\color{red}  #1} }
\def\FS#1{ { {\physical{\bm {#1}} } } }

\def\frameXX{ \physical{e} }
\def\frameXY{ \physical{e} }
\def\frameYX{ \physical{e} }
\def\frameYY{ \physical{e} }

\def\CXXX{ \physical{C} }
\def\CXXY{ \physical{C} }
\def\CXYY{ \physical{C} }
\def\CYYY{ \physical{C} }

\def\grinoXX{ \physical{\psi} }
\def\grinoXY{ \physical{\psi} }
\def\grinoYX{ \physical{\psi} }
\def\grinoYY{ \physical{\psi} }

\def\spinorYY{ \physical{\chi} }
\def\spinorYYY{ \physical{\chi} }

\def\d{ \auxiliary{d_X} }
\def\dX{ \auxiliary{d} }
\def\dY{ \auxiliary{d} }
\def\dYY{ \auxiliary{d} }

\def\f{\auxiliary{f} }

\def\l{ \auxiliary{\lambda} }
\def\y{ \auxiliary{y} }
\def\t{ \auxiliary{t} }
\def\r{ \auxiliary{\rho} }

\def\surfeit{ \compensating{\eta} }
\def\compX{ \compensating{G} }
\def\compY{ \compensating{H} }

\def\calibration{ \physical{F} }

\def\X{X}
\def\T{T}

\def\gravitonXX{\physical{h} }
\def\gravitonXY{\physical{h} }
\def\gravitonYY{\physical{h} }

\def\Ga{\auxiliary{\bm \Gamma}}
\def\Oh{\physical{ {\bm \Omega} }}
\def\I{\bm I}
\def\F{\bm F}
\def\H{\bm H}
\def\Weyl{\physical W}
\def\KK{\physical {\mathcal F}}
\def\Scalar{\physical R}
%\def\F{\bm G_{[2,2] }

\def\G{ \physical{ {\bm G} } }
\def\Torsion{ \physical{ {\bm T} } }
\def\Curvature{ \physical{ {\bm R} } }


\def\TorsionX{ \physical{ {\bm T} } }
\def\TorsionY{ \physical{ {\bm T} } }
%Refinements
\def\TorsionXXX{ \physical{ {\bm T} } }
\def\TorsionXYX{ \physical{ {\bm T} } }
\def\TorsionYYX{ \physical{ {\bm T} } }
\def\TorsionXXY{ \physical{ {\bm T} } }
\def\TorsionXYY{ \physical{ {\bm T} } }
\def\TorsionYYY{ \physical{ {\bm T} } }

%4form Field strength
\def\FSXXXX{\hyperlink{E:FSXXXX} {\FS{G}} }
\def\FSXXXY{\hyperlink{E:FSXXXY} {\FS{G}} }
\def\FSXXYY{\hyperlink{E:FSXXYY} {\FS{G}} }
\def\FSXYYY{\hyperlink{E:FSXYYY} {\FS{G}} }
\def\FSYYYY{\hyperlink{E:FSYYYY} {\FS{G}} }
\def\Torsion{ {\FS{T}} }
\def\Curvature{ {\FS{R}} }

\def\Ztorsion{ { \auxiliary{Z} } }
%%%%%%%%%%%%%%%%%%%%%%%%%%%%%%%%%%%%%%%%%%%%%%%%%%%%%%%%%%%%%%%%%
%%%%%%%%%%%%%%%%%%%%%%%%%%%%%%%%%%%%%%%%%%%%%%%%%%%%%%%%%%%%%%%%%
%%%%%%%%%%%%%%%%%%%%%%%%%%%%%%%%%%%%%%%%%%%%%%%%%%%%%%%%%%%%%%%%%


%
%%%
%%%%%
%%%%%%%%%
%%%%%%%%%%%%%%%%%
%%%%%%%%%%%%%%%%%%%%%%%%%%%%%%%%%
%%%%%%%%%%%%%%%%%%%%%%%%%%%%%%%%%%%%%%%%%%%%%%%%%%%%%%%%%%%%%%%%%
%
%UMDEPP-014-026
%\title{Components of Eleven-dimensional Supergravity from the Superspace Chern-Simons Action}
%%%%%%%%%%%%%%%%%%%%%%%%%%%%%%%%%%%%%%%%%%%%%%%%%%%%%%%%%%%%%%%%
\begin{document}
%\nopagecolor
%\pagecolor{yellow!30!orange}
%\pagecolor{white!60!black}
%\maketitle
%\default@color{white}


%%\begin{center}
%\addtolength{\baselineskip}{.5mm}
%\thispagestyle{empty}
%\begin{flushright}
%{\color{red}\today \\
%{\sc MI-TH-16??}\\
%}
%\end{flushright}

\vspace{10mm}
\begin{center}
%\doublespacing{
\Large \textbf {
Eleven-dimensional Supergravity with Four Off-shell Supersymmetries
}
%}
\end{center} 
\vspace{5mm}
\begin{center} 
%\\[10mm] 
%\\[15mm]
{
William D.\ Linch III
}
\\[5mm]
%~\\[5mm]
{
{\em George P. and Cynthia W. Mitchell Institute for\\
Fundamental Physics and Astronomy, \\
Texas A\&{}M University.}\\
~\\
Caltech Theory Seminar\\
7 August, 2020.\\
~\\
Based on work 
%\cite{Becker:2016xgv, %ATH
%Becker:2016rku,	%NATH
%Becker:2016edk,	%G2Potential
%Becker:2017njd, 	%All CS
%Becker:2017zwe,	%Linearized
%Becker:2018phr}	%Current
with \\
~\\
K.\ Becker, 
M.\ Becker, 
D.\ Butter, 
S.\ Guha, 
S.\ Randall, 
D.\ Robbins,
and
A.\ Sengupta.
}
\end{center} 

\vspace{10pt}


\begin{abstract}
Off-shell supersymmetry of higher-derivative corrections to 2-derivative actions generally leads to an infinite set of ever-higher-derivative corrections. In some cases, the result is a potentially large part of an effective action for some more fundamental theory. While it is generally expected to be impossible to describe more than eight {\color{red} twelve?} Poincar\'e supersymmetries this way, some constructions may be simplified considerably even if only a fraction of them closes off-shell. In this talk, I will describe eleven-dimensional supergravity with four of the 32 supersymmetries manifest and the status of a project to construct the famous $\int C\wedge X_8$ invariant.
\end{abstract}

%%email addresses
%%\vspace*{.5cm}
%\begin{flushleft}
%~\\
%{${}^{\small \mbox\Pisces}$ \href{mailto:wdlinch3@gmail.com}{wdlinch3@gmail.com}}\\
%%\makebox[0pt][t]
%%{$^\text{\Scorpio}$ \href{mailto:siegel@insti.physics.sunysb.edu}{siegel@insti.physics.sunysb.edu}}
%%\makebox[0pt][t]
%\end{flushleft}
\thispagestyle{empty}
\newpage
%\thispagestyle{empty}
\tableofcontents

%\newpage
\setcounter{page}1
%
%%%%%%%%%%%%%%%%%%%%%%%%%%%%%%%%%%%%%%%%%%%%%%%%%%%%%%%%%%%%%%%%
%%%%%%%%%%%%%%%%%%%%%%%%%%%%%%%%%%%%%%%%%%%%%%%%%%%%%%%%%%%%%%%%
\section{Motivation and Strategy}

``M-theory'' is still poorly understood.
It should have (at least) 2-branes and 5-branes (not necessarily fundamental). These have comparable tensions (not parametrically-separated), so a world-volume-like description has not been found.\\
The alternative is to exploit the spacetime properties (\eg symmetries and anomaly structure) together with known limits in the moduli space of theories (\ie string theories). 

We will rely on the following properties of M:
\begin{enumerate}
\item 11D super-Poincar\'e symmetry represented by the minimal multiplet as explained by Cremmer, Julia and Scherk (1978): 11D SG represented in components by a metric $g_{\bm {mn}}$, a gravitino $\psi_{\bm m}^{\bm \alpha}$, and a gauge 3-form $C_{\bm {mnp}}$. 
\item Partial knowledge of the leading higher-derivative corrections to the 11D SG: $\int C \wedge X_8$, with $X_8 \sim R^4+\cdots$ a specific contraction of the fourth exterior power of the curvature 2-form with corrections involving fewer curvatures.
\item The 4-form field strength $G=dC$ does not have integer periods on a generic spinnable Lorenzian 11-manifold. What has integral periods is instead $G - \tfrac14 p_1$ where $p_1 \sim \mathrm{tr} \, R\wedge R$ is the first Pontryagin class (Witten 1997).
\item The long-standing conjecture that this is the leading Wilsonian effective action of M-theory. Schematically,
{%\color{white}
\begin{align}
\label{E:MLEEA}
\boxed{R + R^4 + \textrm{11D, $N=1$ supersymmetry} \Rightarrow \textrm{Leading M-theory LEEA?}}
\tag{$\ast$}
%\nonumber
\end{align}}
\end{enumerate}

Now suppose there were an off-shell, 11D superspace with all 32 supercharges. \\
Then a description of the leading terms in the M-theory LEEA in this superspace must actually contain all orders, at least when all the auxiliary fields are integrated out. \\
The catch is that this superspace does not exist. (For rigid supersymmetry Berkovits (1993) explains that the algebra will not close unless the number of supercharges is 9 or less due to non-associativity of $\mathbb O$.)

Can we make do with fewer than 32 supercharges, say, only 4?\\
In this 11D, $N=1/8$ superspace, the 4 supersymmetries would close on 4 translations giving a 4D, $N=1$ {\em off-shell} subalgebra of the full 11D, $N=1$ on-shell supersymmetry algebra.\\
We would then have to impose the remaining symmetries ``by hand''. \\
A good analogy was provided by Bagger and Galperin \cite{Bagger:1996wp}: $F^2 + F^4 +$ 10D, $N=1$ supersymmetry $\Rightarrow$ Born-Infeld.



%%%%%%%%%%%%%%%%%%%%%%%%%%%%%%%%%%%%%%%%%%%%%%%%%%%%%%%%%%%%%%%%%
%\subsection{Noether Procedure}
%We may represent the situation as follows: The M-theory effective action is of the form $S=S_2 + S_8 +\cdots$ where the subscript denotes the number of derivatives on the bosonic fields. \\
%Then under the lowest-order supersymmetry transformation $\delta_0$, $S_2$ is invariant, but $S_8$ is not. \\
%Supersymmetry of the effective action requires a correction $\delta= \delta_0+\delta_6 +\cdots$ so that $\delta_6 S_2$ cancels the transformation $\delta_0 S_8$. \\
%But then $\delta_6 S_8\neq 0$ so that we require an even-higher-order term in the effective action that may cancel this, and so on.
%This is the usual component Noether procedure--it has been carried out explicitly only to very low orders and even then only partially.\\
%Putting this enormous technical hurdle aside, summing up this infinite number of steps would give a supersymmetric, all-orders-in-derivatives effective action. \\
%It is a long-standing conjecture that this is the leading Wilsonian effective action of M-theory. Schematically,
%\begin{align}
%\label{E:MLEEA}
%R + R^4 + \textrm{11D, $N=1$ supersymmetry} \Rightarrow \textrm{Leading M-theory LEEA?}
%\tag{$\ast$}
%%\nonumber
%\end{align}

%%%%%%%%%%%%%%%%%%%%%%%%%%%%%%%%%
%\paragraph{Born-Infeld}
%
%The situation described above has a spin-1 analog discovered by Born and Infeld (1933). 
%They deformed the Maxwell Lagrangian by higher-derivative corrections in such a way that the field strength would be bounded above: $L\sim \tfrac1{\alpha^2 }[1-\sqrt{1-\alpha^2 F^2}]\approx \tfrac12 F^2 + o(\alpha^2 )$.\\
%It was discovered by Tseytlin {\em half a century} later (1985) that this action is the exact leading LEEA of open (super)string field theory. \\
%A decade later, it was realized by Bagger and Galperin (1997) that this action is fixed by extended supersymmetry, so that we have the analog 
%\begin{align}
%F^2 + F^4 + \textrm{10D, $N=1$ supersymmetry} \Rightarrow \textrm{Born-Infeld}
%\nonumber
%\end{align}
%of the conjecture \eqref{E:MLEEA} above. 

\begin{enumerate}
\item We will attempt to take 4 real supersymmetries off-shell. This is the maximal amount possible with finitely many fields. These symmetries will be taken to close on 4 of the 11 translations.
\item Split the indices of the component fields of the eleven-dimensional super-Poincar\'e multiplet into $4+7$ and embed each one into supermultiplets according to their na\"ive 4D interpretation. 
The M-theory 3-form gives rise to a double complex of differential superforms (non-abelian tensor hierarchy) in this 11D, $N=1/8$ superspace \cite{Becker:2016xgv, Becker:2016rku, Becker:2017njd}.\footnote{Mathematically, this is a {\it Diff}(7)-equivariant chain complex of differential forms on $Y$ with values in a complex of (sheaves of) representations of the 4D, $N=1$ super-Poincar\'e algebra \cite{Becker:2016rku}.} 
\item Construct the 2-derivative action as a covariantization of the associated Chern-Simons invariant, and compute its field equations \cite{Becker:2016edk, Becker:2017zwe, Becker:2017njd, Becker:2018phr}. This tells us how to compare to on-shell theories and is the starting point for studying deformations. 
\item Act repeatedly on the unconstrained superfields in all possible ways with super-translations and consider all possible linear combinations. This gives a large differential algebra containing connections and tensors of our reduced structure algebra. Identifying ``differential ideal'' is tantamount to finding the (reduced) supergeometry.\\
%\xymatrix{
%U \ar[d]_{\bar D} \ar[r]^{D} & DU\ar[d]_{\bar D} \ar[r]^{D} & D^2 U \ar[d]_{\bar D} \ar[r]^{D} &0&&&&\\
%\bar DU\ar[d]_{\bar D} \ar[r]^{D} & [D , \bar D] U \ar[d]_{\bar D} \ar[r]^{D} & \bar D D^2 U\ar[d]_{\bar D} \ar[r]^{D} &0\\
%\bar D^2 U\ar[d]_{\bar D} \ar[r]^{D} & \bar D^2 U \ar[d]_{\bar D} \ar[r]^{D} & \{\bar D^2, D^2\} U\ar[d]_{\bar D} \ar[r]^{D} &0\\
%0&0&0
%}

{%\color{white}
\xymatrix{
\compensating {V} \ar[d]_{\bar D} \ar[r]^{D} & \gauge {DV}\ar[d]_{\bar D} \ar[r]^{D} & \compensating {D^2 V} \ar[d]_{\bar D} \ar[r]^{D} &0&&&&\\
\gauge{\bar DV} \ar[d]_{\bar D} \ar[r]^{D} & \gauge {A_a} \ar[d]_{\bar D} \ar[r]^{D} & \physical{\bar W^\alpha} \ar[d]_{\bar D} \ar[r]^{D} &0\\
\compensating{\bar D^2 V}\ar[d]_{\bar D} \ar[r]^{D} & \physical{W^\alpha} \ar[d]_{\bar D} \ar[r]^{D} & \physical{F_{ab}} \oplus \auxiliary {d} \ar[d]_{\bar D} \ar[r]^{D} &0\\
0&0&0
}
}

{\color{white}
\xymatrix{
\compensating {U} \ar[d]_{\bar D} \ar[r]^{D} & \compensating {DU}\ar[d]_{\bar D} \ar[r]^{D} & \compensating {D^2 U} \ar[d]_{\bar D} \ar[r]^{D} &0&&&&\\
\compensating{\bar DU} \ar[d]_{\bar D} \ar[r]^{D} & \gauge {A_a} \ar[d]_{\bar D} \ar[r]^{D} & \physical{\bar W^\alpha} \ar[d]_{\bar D} \ar[r]^{D} &0\\
\compensating{\bar D^2 U}\ar[d]_{\bar D} \ar[r]^{D} & \physical{W^\alpha} \ar[d]_{\bar D} \ar[r]^{D} & \physical{F_{ab}} \oplus \auxiliary {d} \ar[d]_{\bar D} \ar[r]^{D} &0\\
0&0&0
}
}


\item Construct the full supergeometry of this 11D, $N=1/8$ superspace so that is agrees with the 11D, $N=1$ on-shell supergeometry in the appropriate limit.
This requires finding the 
	\begin{itemize}
	\item field strength super-4-form $\G_{\bm{ABCD}}$ 
	\item torsion 2-form $\Torsion_{\bm{AB}}{}^{\bm C}$ 
	\item curvature 2-form $\Curvature_{\bm{AB} }{}^{\bm{cd} }$
	\end{itemize}
where $\bm A = \bm a, \alpha, \dt \alpha$ and $\bm a = a, i$. 
\item Use the resulting curvature super-2-form tensor to define the  super-4-form $\Curvature^2$ and super-8-form $\Curvature^4$, and use them to deform the 2-derivative action.
%\item Deform the action by the super-8-form and perform the Witten shift on the super-4-form field strength $\G=dC$. 
\end{enumerate}







%%%%%%%%%%%%%%%%%%%%%%%%%%%%%%%%%%%%%%%%%%%%%%%%%%%%%%%%%%%%%%%%
%%%%%%%%%%%%%%%%%%%%%%%%%%%%%%%%%%%%%%%%%%%%%%%%%%%%%%%%%%%%%%%%
\section{Putting it up to 11(D SG in $N=1/8$ Superspace)}
Following the strategy for spin-1, we first decompose 11D SG fields 
\begin{align}
\begin{array}{cccclcc}
	e_{\bm m}{}^{\bm a} &\to &
		e_{m}{}^{a} ~,~ 
		e_{i}{}^{a} ~,~ 
		e_{m}{}^{j} ~,~ 
		{\color {teal} e_i{}^j }
&\with & a, m = 0,1,2,3
\\
	\psi_{\bm m}^{\bm \alpha}&\to &
		\psi_{m}^{\alpha} ~,~ 
		\psi_{m}^{\alpha i} ~,~ 
		\psi_{j}^{\alpha } ~,~ 
		\psi_{j}^{\alpha i} ~,~
&\and & i,j = 1,\dots, 7	
\\ 
	C_{\bm {mnp}} &\to &C_{mnp}~,~ C_{mn \, i} ~,~ C_{m \, ij}~,~ C_{ijk}
&\and & \alpha = 1,2	
\end{array}
\nonumber
\end{align}
and embed them into superfields. 
This requires the following surprising fact about Riemannian 7-manifolds:
For any orientable and spinnable such smooth manifold, 
%\begin{theorem}
%Let $(Y, g)$ be a smooth Riemannian manifold that is oriented ($w_1=0$) and spin ($w_2$=0). 
%Then 
%\end{theorem}
there is a 3-form $\varphi$ such that 
{%\color{white}
\begin{align}
{\color {teal} \sqrt{g}g_{ij} } = -\tfrac1{144}\varepsilon^{k_1\cdots k_7} \phi_{i k_1k_2} \phi_{j k_3 k_4} \phi_{k_5k_6k_7}
\nonumber 
\end{align}
}
This 3-form is then embedded into a chiral superfield $\Phi_{ijk}$ together with the 35 pseudoscalars. 
All the other bosons get their own superfields by the na\"ive assignment \cite{Becker:2017zwe}
{%\color{white}
\begin{align}
U^a &\sim \cdots 
	+ (\theta \sigma^m \bar \theta) \frameXX_m{}^a 
	+ \theta^2 (\sigma^m \bar \theta)_\alpha \grinoXX_m{}^\alpha 
	+ \bar \theta^2 (\theta \sigma^m )_{\dt \alpha} \bar \grinoXX_m{}^{\dt \alpha }
	+ \theta^2 \bar \theta^2 \dX^a
\cr
\Psi^{\alpha i} &\sim \cdots 
	+ (\theta \sigma^m \bar \theta) \grinoXY_m{}^{\alpha i}
	+ \theta^2 (\bar \theta \bar \sigma_a)^\alpha \y_i^a
	+ \bar \theta^2 (\theta \sigma^{ab})^\alpha \t_{ab i} 
	+ \theta^2 \bar \theta^2 \r^i_\alpha
\cr
\mathcal V^i &\sim \cdots 
	+ (\theta \sigma^m \bar \theta) \frameXY_m^i 
	+ \bar \theta^2 \theta^\alpha \surfeit_\alpha^i
	+ \theta^2 \bar \theta_{\dt \alpha}\bar \surfeit^{\dt \alpha i}
	+ \theta^2\bar \theta^2 \dY^i
\cr
X &\sim \cdots 
	+ \theta^2 \compX
	+ \bar \theta^2 \bar \compX
	+ \theta \sigma_m \bar \theta \mathcal \, \varepsilon^{mnpq} \CXXX_{npq}
	+ \bar \theta^2 \theta^\alpha \surfeit_\alpha
	+ \theta^2 \bar \theta_{\dt \alpha}\bar \surfeit^{\dt \alpha}
	+ \theta^2\bar \theta^2 \d
\cr
\Sigma_i^\alpha &\sim \cdots	
	+ [\theta^\alpha \compY_i + (\theta \sigma^{mn})^\alpha \CXXY_{mn \, i} ] 
	+  \theta^2 \tilde \surfeit_i^\alpha
\cr
V_{ij} &\sim \cdots 
	+  (\theta \sigma^m \bar \theta)  \CXYY_{m\, ij}
	+ \bar \theta^2 \theta^\alpha \spinorYY_{\alpha \, ij} 
	+ \theta^2 \bar \theta_{\dt \alpha}\bar \spinorYY_{ij}^{\dt \alpha}
	+ \theta^2\bar \theta^2 \dYY_{ij}
\cr
\Phi_{ijk} &\sim \CYYY_{ijk} + i \calibration_{ijk}  
	+ \theta^\alpha \spinorYYY_{\alpha\, ijk}
	+ \theta^2 \f_{ijk}
\nonumber
\end{align}
}
Notes:
\begin{itemize} 
	\item {\color{PineGreen} Gauge fields} transform canonically. Their (double) curls are field strength, torsion, and curvature 2-form components. 
	\item {\color{Orange} Auxiliary fields} are invariant. 
	\item {\color{Red} Compensators} suffer St\"uckelberg shifts under ``pre-gauge'' transformations. They do not appear in the spectrum of gauge(-invariant) superfields.
\end{itemize}


The spins $\leq 1$ fields in the list form a ``non-abelian tensor hierarchy'' of superfield.
Such homological objects have a complex of field strengths:
%\begin{align}
%E_{ijkl} &= 4\partial_{[i} \Phi_{jkl]}
%	\cr
%F_{ijk} &= \tfrac1{2i}\left( \Phi_{ijk} - \bar \Phi_{ijk}\right) - 3\partial_{[i} V_{jk]}
%	\cr
%W_{\alpha ij} &= -\tfrac14 \bar {\mathcal D}^2 {\mathcal D}_\alpha V_{ij}
%	+2\partial_{[i} \Sigma_{\alpha j]}
%	+ \mathcal W_\alpha^k \Phi_{ijk}
%	\cr
%H_i &= \tfrac1{2i}\left({\mathcal D}^\alpha \Sigma_{\alpha i} - \bar {\mathcal D}_{\dt \alpha} \bar \Sigma_i^{\dt \alpha} \right)
%	-\partial_i X 
%	+\left[  \mathcal W^{\alpha j} {\mathcal D}_\alpha  V_{ij}
%		+\tfrac12 {\mathcal D}^\alpha \mathcal W_\alpha^j   V_{ij}
%	+\hc \right]
%	\cr
%G&= -\tfrac14 \bar {\mathcal D}^2 X 
%	+ \mathcal W^{\alpha i} \Sigma_{\alpha	i}
%\nonumber
%\end{align}
{%\color{white}
\begin{align}
E_{ijkl} &= 4\partial_{[i} \Phi_{jkl]}
	\cr
F_{ijk} &= \tfrac1{2i}\left( \Phi_{ijk} - \bar \Phi_{ijk}\right) - 3\partial_{[i} V_{jk]}
	\cr
W_{\alpha ij} &= -\tfrac14 \bar {\mathcal D}^2 {\mathcal D}_\alpha V_{ij}
	+2\partial_{[i} \Sigma_{\alpha j]}
	+ \mathcal W_\alpha^k \Phi_{ijk}
	\cr
H_i &= \tfrac1{2i}\left({\mathcal D}^\alpha \Sigma_{\alpha i} - \bar {\mathcal D}_{\dt \alpha} \bar \Sigma_i^{\dt \alpha} \right)
	-\partial_i X 
	+\left[  \mathcal W^{\alpha j} {\mathcal D}_\alpha  V_{ij}
		+\tfrac12 {\mathcal D}^\alpha \mathcal W_\alpha^j   V_{ij}
	+\hc \right]
	\cr
G&= -\tfrac14 \bar {\mathcal D}^2 X 
	+ \mathcal W^{\alpha i} \Sigma_{\alpha	i}
\nonumber
\end{align}
}
It is gauged by the non-abelian KK vector multiplet $\mathcal V^i$ with field gauge-chiral strength $\mathcal W_\alpha^i$.

In addition there are some combinations of the other three superfields $\{U^a, \Psi_i, \mathcal V^i\}$ that appear (mixed with a little $X$)
{%\color{white}
\begin{align}
S &:=\tfrac12(G+\bar G) -\tfrac14 [D \sigma_a \bar D] U^a
\cr
P &:=\tfrac 1{2i} (G-\bar G) + \partial_a U^a
\cr
\X_i^a &:= \tfrac i4 ( \bar D \bar \sigma^a \Psi_i+ D\sigma^a\bar \Psi_i ) - \partial_i U^a
\cr
\T_i &:= \tfrac1{2i} ( D\Psi_{i} - \bar D \bar \Psi_i) + H_i
%\cr
%\bm \Gamma^{\un a}_i &:= -i\bar D^2 \X^{\un a}_i
%	- i \bar D^{\dt \alpha} D^\alpha \T_i 
\nonumber
\end{align}
}

%%%%%%%%%%%%%%%%%%%%%%%%%%%%%%%%%%%%%%%%%%%%%%%%%%%%%%%%%%%%%%%%
\subsection{``Components''}
The formul\ae{} above can be partially inverted for the (linearized) components.\footnote{Need to define $\tau, Z$.} 
%\begin{align}
%\CXXX_{mnp}&:=- \tfrac12\varepsilon_{mnpq}\left[ 
%	\tfrac12 (\bar \sigma^q)^{\dt \delta \delta} [D_\delta, \bar D_{\dt \delta}] X 
%	+(D^2 +\bar D^2) U^d
%	\right]
%\cr
%\CXXY_{mn \, i}  &:=\tfrac12 D_{(\alpha} \left[ \Sigma_{\beta) i} + \Psi_{\beta) i}\right]
%\cr
%\CXYY_{m\, ij}&:=\tfrac 12 [D_\alpha, \bar D_{\dt \alpha}] V_{ij} + \varphi_{ijk} \X_{\un a}^k
%\cr
%\CYYY_{ijk} &:=\tfrac12\left( \Phi_{ijk} + \bar \Phi_{ijk}\right)
%	+\tfrac12\psi_{ijkl} \T^l
%\nonumber
%\end{align}
%%%%%%%%%%%%%
%%%% Graviton
%%%%%%%%%%%%%
%\begin{align}
%\frameXX_m{}^a &:=-\tfrac12 [D_\alpha, \bar D_{\dt \alpha} ] \, U^{\un b} 	-\tfrac23 \delta_\alpha^\beta \delta_{\dt \alpha}^{\dt \beta} S
%\cr
%\frameXY_m{}^i &:=\tfrac12 [D_\alpha, \bar D_{\dt \alpha} ] \mathcal V^i 
%\cr
%\frameYX_i{}^b &:=
%	- \tfrac12\left( \bar D_{\dt \alpha} \Psi_{\alpha i} - D_{\alpha} \bar \Psi_{\dt \alpha i}\right)
%\cr
%\frameYY_i{}^j&:= \tfrac14 \varphi^{j kl} F_{ikl} 
%	- \tfrac1{36} \delta_i^j \varphi^{klm} F_{klm}
%\nonumber
%\end{align}
{%\color{white}
\begin{align}
-4 \CXXX_{mnp}&:= \varepsilon_{mnpq}\left(
	 [D\sigma^q\bar D] X
%(\bar \sigma^q)^{\dt \delta \delta} [D_\delta, \bar D_{\dt \delta}] X 
	+2(D^2 +\bar D^2) U^d
	\right)
\cr
2\CXXY_{mn \, i}  &:= (D\sigma_{mn}\Sigma_{i} ) 
	+ (D\sigma_{mn}\Psi_{ i})
\cr
-4 \CXYY_{m\, ij}&:= 
%(\bar \sigma_m)^{\dt \alpha \alpha} [D_\alpha, \bar D_{\dt \alpha}] 
	[D\sigma_m \bar D] V_{ij} -4 \varphi_{ijk} \X_{a}^k
\cr
2 \CYYY_{ijk} &:= \Phi_{ijk} + \bar \Phi_{ijk} +\psi_{ijkl} \T^l
\nonumber
\end{align}
%%%%%%%%%%%%
%%% Graviton
%%%%%%%%%%%%
\begin{align}
4\frameXX_m{}^a &:=
%(\bar \sigma_m)^{\dt \alpha \alpha} [D_\alpha, \bar D_{\dt \alpha} ] \, 
	[D\sigma_m \bar D]U^{a} 
	-\tfrac83  \delta_m^a S
\cr
-4 \frameXY_m{}^i &:= 
%(\bar \sigma_m)^{\dt \alpha \alpha} [D_\alpha, \bar D_{\dt \alpha} ] 
	[D\sigma_m \bar D]\mathcal V^i 
\cr
4 \frameYX_i{}^b &:=
	 \bar D \sigma^b \Psi_{i} - D \sigma^b \bar \Psi_{ i}
\cr
4\frameYY_i{}^j&:=  \varphi^{j kl} F_{ikl} 
	- \tfrac1{9} \delta_i^j \varphi^{klm} F_{klm}
\nonumber
\end{align}
%%%%%%%%%%%%%
%%%% Gravitino
%%%%%%%%%%%%%
%\begin{align}
%\grinoXX_m{}^\alpha &:=-\tfrac{i}8 \bar D^2 D^{\beta} U_{\alpha\dt \alpha} 
%	+\tfrac i{3} \delta_\alpha^\beta
%		\bar D_{\dt \alpha}S
%\cr
%\grinoYX_i{}^\beta &:= \tfrac i8 \bar D^2 \Psi_i^\alpha
%	+\tfrac1{12}\left[
%	2 \bar D_{\dt \alpha} \X_i^{~\un a} - D^\alpha \T_i
%	-4i \mathcal W_i^\alpha
%	\right]
%\cr	
%\grinoXY_m{}^{\alpha i}&:=-\tfrac i2 [D_\alpha , \bar D_{\dt \alpha}] \Psi^{\beta j}
%	-2 \partial^j D_\alpha U^\beta{}_{\dt \alpha}
%	- \delta_\alpha^\beta  \bar D_{\dt \alpha} \T^j
%	-4 \delta_\alpha^\beta \bar \grinoXY^j{}_{\dt \alpha}
%\cr
%\grinoYY_i{}^{\beta j} &:= i h_i{}^j(\spinorYYY^\beta)
%	+\partial_{i}\Psi^{\beta j}
%	-\tfrac i{3} \varphi_i{}^{jk}\left[ 
%		\tfrac1{12}\psi_k{}^{lmn}\spinorYYY^\beta_{lmn}
%		+i D^\beta \T_k\right]
%	- \pi_{\bm{14}}{}_i{}^j  
%	( \widehat W^\beta)
%\nonumber
%\end{align}
%%%%%%%%%%%%%
%%%% Alternative combinations 
%%%%%%%%%%%%%
%\begin{align}
%\spinorYY_{ij}^\alpha &:=\widehat W_{\alpha ij} 
%	+i\varphi_{ijk} (\mathcal W_\alpha^k -\tfrac14 \bar D^2 \Psi_\alpha^k)
%\cr
%\spinorYYY_{ijk}^\alpha&:= D^\alpha \left( F_{ijk} - 3i \varphi_{l[ij} \partial_{k]} \mathcal V^l\right)
%\nonumber
%\end{align}
%%%%%%%%%%%%%
%%%% Auxiliary fields
%%%%%%%%%%%%%
%\begin{align}
%\dX^a &:=\tfrac18 D^\beta \bar D^2 D_\beta U_{\un a} 
%	+ \tfrac16 [ D_\alpha, \bar D_{\dt \alpha}] S
%	+\partial_{\un a} P
%\cr
%\d&:= -\tfrac14 (D^2 + \bar D^2) S
%\cr
%3\dY^i&:=D^\alpha \mathcal W_\alpha^i 
%		-\tfrac18\left(
%			D^\alpha \bar D^2 \Psi_\alpha^i 
%			+\bar D_{\dt \alpha} D^2 \bar \Psi^{\dt \alpha i}
%		\right)
%		+2\partial^i P
%\cr
%4\dYY_{ij}&:=	\tfrac12 \left( 
%		D^\alpha \spinorYY_{\alpha \, ij}
%		+\bar D_{\dt \alpha} \bar \spinorYY_{ij}^{\dt \alpha }
%	\right)
%	-\tfrac43 \varphi_{ijk} \partial^k S
%		-4 \varphi_{ijk} (\tau_1)^k 
%		-2 (\tau_2)_{ij}
%\cr
%\f_{ijk}&:=-18 i \bar D^2 \left[ G^{ijk, lmn} \bar Z_{lmn}\right]
%	- \tfrac 13 \epsilon^{ijk lmnp}\partial_l \bar Z_{mnp}
%\cr
%\y_i^a&:= - \bar D_{\dt \alpha} \l_\alpha^i
%\cr
%2 \t_{ab i} &:=D_{(\alpha} \mathcal W_{\beta) }^k
%	-\tfrac i2 \left[  3 \bar D^{\dt \alpha} D_{(\beta}   
%	+D_{(\beta} \bar D^{\dt \alpha}
%	\right] \X_{\alpha) \dt \alpha}^k 
%\cr
%\r^i_\alpha&:=-\tfrac14 D^\beta \bar D^2 D_{(\alpha} \Psi_{\beta)}^i 
%	- i \partial^i \left[
%	D_\alpha G +\tfrac12 \bar D^{\dt \alpha} D^2 U_{\un a}
%	\right]
%\nonumber
%\end{align}
%%%%%%%%%%%%
%%% Gravitino
%%%%%%%%%%%%
\begin{align}
8i \grinoXX_m{}^\beta &:= \bar D^2 D^{\beta} U_{m} 
	+\tfrac 4{3} (\sigma_m\bar D)^\beta S
\cr
-8i \grinoYX_i{}^\beta &:= \bar D^2 \Psi_i^\beta
	-\tfrac {2i}{3}\left[
	2 (\sigma_m\bar D)^\beta \X_i^{m} - D^\alpha \T_i
	-4i \mathcal W_i^\alpha
	\right]
\cr	
-4i\grinoXY_m{}^{\beta j}&:=
%(\bar \sigma_m)^{\dt \alpha \alpha} [D_\alpha , \bar D_{\dt \alpha}] 	
	[D\sigma_m \bar D]\Psi^{\beta j}
	-4i \partial^j (D\bar \sigma_m \sigma_n )^\beta U^n
	-2i (\sigma_m\bar D )^\beta \T^j
	-8i (\sigma_m \bar \grinoXY^j)^\beta 
\cr
\grinoYY_i{}^{\beta j} &:= i h_i{}^j(\spinorYYY_{ijk}^\beta)
	+\partial_{i}\Psi^{\beta j}
	-\tfrac i{3} \varphi_i{}^{jk}\left[ 
		\tfrac1{12}\psi_k{}^{lmn}\spinorYYY^\beta_{lmn}
		+i D^\beta \T_k\right]
	- \pi_{\bm{14}}
	( \spinorYY_{i}^{\beta\, j})
\nonumber
\end{align}
%%%%%%%%%%%%
%%% Alternative combinations 
%%%%%%%%%%%%
\begin{align}
\spinorYY_{ij}^\alpha &:=\widehat W_{ij}^\alpha 
	+i\varphi_{ijk} (\mathcal W^{\alpha k} -\tfrac14 \bar D^2 \Psi^{\alpha k})
\cr
\spinorYYY_{ijk}^\alpha&:= D^\alpha \left( F_{ijk} - 3i \varphi_{l[ij} \partial_{k]} \mathcal V^l\right)
\nonumber
\end{align}
%%%%%%%%%%%%
%%% Auxiliary fields
%%%%%%%%%%%%
\begin{align}
8 \dX^a &:= D^\beta \bar D^2 D_\beta U^a
	- \tfrac 2{3} 
%(\bar \sigma^a)^{\dt \alpha \alpha} [ D_\alpha, \bar D_{\dt \alpha}] 
	[D\sigma^a \bar D]S
	+ 8 \partial^{a} P
\cr
-4 \d&:=  (D^2 + \bar D^2) S
\cr
3\dY^i&:=D^\alpha \mathcal W_\alpha^i 
		-\tfrac18\left(
			D^\alpha \bar D^2 \Psi_\alpha^i 
			+\bar D_{\dt \alpha} D^2 \bar \Psi^{\dt \alpha i}
		\right)
		+2\partial^i P
\cr
8\dYY_{ij}&:=
		D^\alpha \spinorYY_{\alpha \, ij}
		+\bar D_{\dt \alpha} \bar \spinorYY_{ij}^{\dt \alpha }
	-\tfrac83 \varphi_{ijk} \partial^k S
		-8 \varphi_{ijk} (\tau_1)^k 
		-4 (\tau_2)_{ij}
\cr
\f_{ijk}&:=-18 i \bar D^2 \left[ G^{ijk, lmn} \bar Z_{lmn}\right]
	- \tfrac 13 \epsilon^{ijk lmnp}\partial_l \bar Z_{mnp}
\cr
2\y_i^a&:=  (\bar D \bar \sigma^a \l_i)
\cr
2 \t_{ab i} &:=(D \sigma_{ab} \mathcal W_{}^k)
	+\tfrac i2 \left[  3 (\bar D^{\dt \alpha} \bar \sigma^c\sigma_{ab}  D)
	+ (D \sigma_{ab} \sigma^c \bar D)
	\right] \X_{c}^k 
\cr
\l^i_\alpha 
%&:= D_\alpha \T^i - 2(\sigma_a \bar D)_\alpha \X_i^a
%	+2i \left[
%	\mathcal W^i_\alpha - \tfrac i2\varphi^{i jk} \widehat W_{\alpha jk}
%	- \tfrac 1{12} \psi^{i jkl} \spinorYYY_{\alpha jkl}
%	\right]
%\cr
	&:= 12 \grinoYX^i{}_\alpha 
	+ \varphi^{ijk} \spinorYY_{\alpha\, jk}
	-\tfrac i6 \psi^{i jkl} \spinorYYY_{\alpha\, jkl}
\cr
-4 \r^i_\alpha&:= D^\beta  \bar D^2 D_{(\alpha} \Psi_{\beta)}^i 
	+4i \partial^i \left[
	D_\alpha G +\tfrac12 \bar D^{\dt \alpha} D^2 U_{\alpha \dt \alpha}
	\right]
\nonumber
\end{align}
Notes:
\begin{itemize} 
\item These are all superfields. 
\item Supersymmetry acts by supertranslations $\delta = -\epsilon D - \bar \epsilon \bar D$.
\item {\color{PineGreen} Gauge fields} transform canonically. Their (double) curls are field strength, torsion, and curvature 2-form components. {\color{Orange} Auxiliary fields} are invariant. 
\item {\color{Red} Compensators} suffer St\"uckelberg shifts under ``pre-gauge'' transformations. They do not appear in the spectrum of gauge(-invariant) superfields.
\item There is only one dimension-1/2 invariant $\l^i_\alpha$. Because of this field, the definition of the gravitino components is ambiguous. We have chosen a minimal prescription in which any such a term is absent.
\end{itemize}

} %end white

%%%%%%%%%%%%%%%%%%%%%%%%%%%%%%%%%%%%%%%%%%%%%%%%%%%%%%%%%%%%%%%%
\subsection{Off-Shell Supersymmetry}

The Wess-Zumino transformations of the 3-form field are 
{%\color{white}
\begin{align}
\delta_\epsilon \CXXX_{abc} &= -6 (\epsilon \sigma_{[ab} \grinoXX_{c]})
	+\mathrm{h.c.}
\cr
\delta_\epsilon \CXXY_{ab\, i} &= - 4i (\epsilon \sigma_{[a} \bar \grinoXY_{b] i}) 
	-2 (\epsilon \sigma_{ab} \grinoYX_{i})
	+\mathrm{h.c.}
\cr
\delta_\epsilon \CXYY_{a\, ij} &= \varphi_{ijk} (\epsilon \grinoXY_{a}{}^k ) 
	-4i\delta_{k[i} (\epsilon \sigma_{a} \bar \grinoYY_{j]}{}^k)
	-\tfrac16 \varphi_{ijk} (\epsilon \sigma_{a} \bar \l_{}^k)
	+\mathrm{h.c.}
\cr
\delta_\epsilon \CYYY_{ijk} &= \varphi_{l[ij} (\epsilon \grinoYY_{k]}{}^l)
	+\mathrm{h.c.}
\nonumber	
\end{align}
}
Note the appearance of the dimension-1/2 auxiliary field $\l_\alpha^i$. 

We now proceed to the off-shell metric transformations. These are unchanged from their on-shell values by our minimal choice of gravitini:
{%\color{white}
\begin{align}
\delta_\epsilon \gravitonXX_{ab} &= -4 (\epsilon \sigma_{(a} \bar \grinoXX_{b)})
	+\mathrm{h.c.}
\cr
\delta_\epsilon \gravitonXY_{a i} &= - i (\epsilon\grinoXY_{a i}) -2(\epsilon \sigma_{a} \bar \grinoYX_{i}) 
	+\mathrm{h.c.}
\cr
\delta_\epsilon \gravitonYY_{ij} &= -2i (\epsilon\grinoYY_{(i, j)})
	+\mathrm{h.c.}
\nonumber	
\end{align}
}
In principle there could have been a contribution by $\l$ to the off-shell $h_{a i}$ transformation, but in contrast to the 3-form transformations, it does not enter (unless we expressly redefine the gravitino components). 
No other contributions are possible because there is there is no other dimension-1/2 invariant with the required $G_2$ representations.
{%\color{white}
\begin{align}
\delta_\epsilon \grinoXX_a^\beta &\simeq
	-i  (\epsilon \eta_{ab} - \epsilon \sigma_a \bar \sigma_b)^\beta \dX^b
	-\tfrac 16 (\bar \sigma_a \bar \epsilon)^{\beta} \d
	+\tfrac 1{72} (\bar \sigma_a \bar \epsilon)^{\beta} \varepsilon^{bcde} \G_{bcde}
\cr	
\delta_\epsilon \grinoYX_i^\beta &\simeq
	-\tfrac12 (\epsilon \sigma_{ab})^\beta \omega_i{}^{ab}
	+\tfrac i{12} (\epsilon \sigma_{ab})^\beta \t_i^{ab}
	-\tfrac i2 \epsilon^\beta \dX_i  
	+ \tfrac i{12} (\bar \epsilon \bar \sigma_a )^\beta \Ga_i^{a}
\cr
\delta_\epsilon \grinoXY_{a}{}^{\beta j} 
	&\simeq 
	\tfrac{5i}{12} (\delta_a^b\epsilon 
		- \epsilon \sigma_a \bar \sigma^b)^\beta\,  \bar \Ga_b^j
	- \epsilon^\beta \left[
		\tilde \omega_a{}^{kl}\varphi_{kl}{}^j 
		+\tfrac i4 \Ga^j_a
		-\tfrac18\left(
			\bar D\bar \sigma_a \l^j - D\sigma_a \bar \l^j 
		\right) 
	\right]
	\cr
	& + (\bar \epsilon \bar \sigma_b)^\beta \left[ 
		- \delta_a^b \d^j - \omega_a{}^{bj} 
		-\tfrac12 \varepsilon_a{}^{bcd} \t^j_{cd}
		\right]
	+ (\bar \epsilon \bar \sigma^{bc}\bar \sigma_a)^\beta \left[ 
		\tfrac i3 \varphi^{jkl} \G_{bc\, kl} 
		+\tfrac16 (\bar D \bar \sigma_{bc} \bar\l^j)
		\right]		
\cr
\delta_\epsilon \grinoYY_{i}{}^{\beta j} 
	&\simeq
	-i(\epsilon \sigma^{ab})^\beta \pi_{\bm 14} \G_{ab\,i}{}^j
	\cr
	&-\tfrac 12 \epsilon^\beta \left[
		2\tilde \omega_i{}^{kl}\varphi_{kl}{}^j
		+\bar \Oh_i{}^j
		-\tfrac16 \varphi_i{}^{jk} \left(
			3i \dY_k 
			+\tfrac23 \varphi_k{}^{lm} \dYY_{lm}
			+\tfrac1{36} \psi_k{}^{lmn} \bar \f_{lmn}
		\right)
		\right]
	\cr
	&+ (\bar \epsilon \bar \sigma^a)^\beta \left[
		2\omega_i{}^{bj}  
		-i\gravitonYY_i^j(\G^b{}_{ijk})
		+\tfrac i6 \varphi_i{}^{jk} \Ga_k^b 
		-\tfrac 1{12} \varphi_i{}^{jk} \bar D \bar \sigma^b \l_k
		\right]
\nonumber	
\end{align}


\begin{align}
\Ga^{a}_i 
%&:= -i\bar D^2 \X^{a}_i+\tfrac i2 \bar D \bar \sigma^a D^\alpha \T_i 
%\cr&
&:=\tfrac i{3!}\varepsilon^{abcd} \G_{bcd \,i}
	+\tfrac16 \psi_{ijkl} \G^{a jkl} 
	+\tfrac1{4i} ( \bar D \bar \sigma^a \l_i + D \sigma^a \bar \l_i )
\cr
\bar \Oh_i{}^j &:= 2\partial_k [ \gravitonYY_{il} - 2i \partial_{(i}\mathcal V_{l)} ]  \varphi^{kl j} 
 	+iD^2 [  \gravitonYY_i{}^j - 2i \partial_{(i}\mathcal V^{j)} ]
	+\tfrac16\partial_i \psi^{jklm} Z_{klm}
	- q_i{}^j	
\nonumber
\end{align}

}



















%\begin{align}
%\frameXX_m{}^a &:=-\tfrac12 [D_\alpha, \bar D_{\dt \alpha} ] \, U^{\un b} 	-\tfrac23 \delta_\alpha^\beta \delta_{\dt \alpha}^{\dt \beta} S
%\cr
%\grinoXX_m{}^\alpha &:=-\tfrac{i}8 \bar D^2 D^{\beta} U_{\alpha\dt \alpha} 
%	+\tfrac i{3} \delta_\alpha^\beta
%		\bar D_{\dt \alpha}S
%\cr
%\dX^a &:=
%\cr
%\grinoXY_m{}^{\alpha i}&:=\tfrac i8 \bar D^2 \Psi_i^\alpha
%	+\tfrac1{12}\left[
%	2 \bar D_{\dt \alpha} \X_i^{~\un a} - D^\alpha \T_i
%	-4i \mathcal W_i^\alpha
%	\right]
%\cr
%\y_i^a&:=
%\cr
%\t_{ab i} &:=
%\cr
%\r^i_\alpha&:=
%\cr
%\frameXY_m{}^i &:=\tfrac12 [D_\alpha, \bar D_{\dt \alpha} ] \mathcal V^i 
%\cr
%\frameYX_i{}^b &:=
%	- \tfrac12\left( \bar D_{\dt \alpha} \Psi_{\alpha i} - D_{\alpha} \bar \Psi_{\dt \alpha i}\right)
%\cr
%\dY^i&:=
%\cr
%\cr
%\CXXX_{mnp}&:=- \tfrac12\varepsilon_{abcd}\left[ 
%	\tfrac12 (\bar \sigma^d)^{\dt \delta \delta} [D_\delta, \bar D_{\dt \delta}] X 
%	+(D^2 +\bar D^2) U^d
%	\right]
%\cr
%\d&:=
%\cr
%\CXXY_{mn \, i} ] &:=\tfrac12 D_{(\alpha} \left[ \Sigma_{\beta) i} + \Psi_{\beta) i}\right]
%\cr
%\CXYY_{m\, ij}&:=\tfrac 12 [D_\alpha, \bar D_{\dt \alpha}] V_{ij} + \varphi_{ijk} \X_{\un a}^k
%\cr
%\spinorYY_{\alpha \, ij} &:=
%\cr
%\grinoYX &:= -\tfrac i2 [D_\alpha , \bar D_{\dt \alpha}] \Psi^{\beta j}
%	-2 \partial^j D_\alpha U^\beta{}_{\dt \alpha}
%%\cr
%	- \delta_\alpha^\beta  \bar D_{\dt \alpha} \T^j
%	-4 \delta_\alpha^\beta \bar \grinoXY^j{}_{\dt \alpha}
%\cr	
%\dYY_{ij}&:=
%\cr
%\CYYY_{ijk} &:=\tfrac12\left( \Phi_{ijk} + \bar \Phi_{ijk}\right)
%	+\tfrac12\psi_{ijkl} \T^l
%\cr
%\grinoYY &:= i h_i{}^j(\spinorYYY^\beta)
%	+\partial_{i}\Psi^{\beta j}
%	-\tfrac i{3} \varphi_i{}^{jk}\left[ 
%		\tfrac1{12}\psi_k{}^{lmn}\spinorYYY^\beta_{lmn}
%		+i D^\beta \T_k\right]
%	- \pi_{\bm{14}}{}_i{}^j  
%	( \widehat W^\beta)
%\cr
%\calibration_{ijk}  &:=
%\cr
%\spinorYYY_{\alpha\, ijk}&:=
%\cr
%\f_{ijk}&:=
%\cr
%\frameYY_i{}^j&:= \tfrac14 \varphi^{j kl} F_{ikl} 
%	- \tfrac1{36} \delta_i^j \varphi^{klm} F_{klm}
%\nonumber
%\end{align}



%%%%%%%%%%%%%%%%%%%%%%%%%%%%%%%%%%%%%%%%%%%%%%%%%%%%%%%%%%%%%%%%
%%%%%%%%%%%%%%%%%%%%%%%%%%%%%%%%%%%%%%%%%%%%%%%%%%%%%%%%%%%%%%%%
\section{Action}



%%%%%%%%%%%%%%%%%%%%%%%%%%%%%%%%%%%%%%%%%%%%%%%%%%%%%%%%%%%%%%%%
\subsection{Chern-Simons-like Action}
They also have a unique Chern-Simons-like invariant \cite{Becker:2016xgv,Becker:2016rku,Becker:2017njd}:
%\begin{align}
%	-12 \kappa^2 \mathcal L_{CS} &= i\int d^2 \theta \, \left[ i \Phi_{[0,3]}\wedge  \mathbb G_{[4,4]}  + \Sigma^\alpha_{[2,1]} \wedge \mathbb W_{[2,6] \alpha }\right] 
%	\cr&
%	+\int d^4 \theta  \left[ V_{[1,2]} \wedge \mathbb H_{[3, 5]} - X_{[3,0]} \mathbb F_{[1,7]}\right]
%	+\hc 
%\nonumber
%\end{align}
{%\color{white}
\begin{align}
-12 L_{CS} &= 
i \int d^2 \theta  \,   \Phi_3 \left[
	{\color{TealBlue} \partial_1\Phi_3 G_0} + {\color{blue} \tfrac12 W_2W_2} - \tfrac i4 \bar {{\mathcal D}}^2 (  F_3 H_1) 
	\right]
\cr&+ \int d^4 \theta  \,  V_2 \left[  
	\partial_1\left( \Phi_3 + \bar \Phi_3 \right) H_1 + \omega(W^1, F_3) 
	- i {\mathcal D} F_3\iota_{\mathcal W^1}F_3 
	+ i \bar {\mathcal D}F_3 \iota_{\bar{\mathcal W}^{1} }F_3 
	\right]
\cr&+i \int d^2 \theta  \, \Sigma_1 \left[
	\partial_1\Phi_3 W_2 -\tfrac i4 \bar {{\mathcal D}}^2 (  F_3 {\mathcal D}F_3) 
	\right]
\cr&- \int d^4 \theta  \, X_0 \left[  
	\partial_1\left( \Phi_3 + \bar \Phi_3 \right) F_3
	\right]
	+\mathrm{h.c.} 	
\nonumber
\end{align}
}
Some observations:
\begin{itemize}
	\item 
		${\color{teal} G \Phi \wedge d \Phi }$ contains the ``$G_2$ superpotential'' (Gukov 1999). 
	\item 
		${\color{blue} \tfrac12 \Phi \wedge W^\alpha \wedge W_\alpha}$ is the holomorphic gauge-kinetic function. 
	\item Every term is required for non-abelian tensor hierarchy gauge invariance: Had we not known this action existed, we would have discovered it by starting from any single particular term.
%	\item Local 4D, N=1 super-conformal invariance.
	\item There is no superpotential. The only F-term integrals of the theory are the ones that appear here.
\end{itemize}


%%%%%%%%%%%%%%%%%%%%%%%%%%%%%%%%%%%%%%%%%%%%%%%%%%%%%%%%%%%%%%%%
\subsection{Supervolume}
The D-term integral
{%\color{white}
\begin{align}
	S_{vol} = -\frac 3{\kappa^2} {\color{Orange} \int d^7 y} {%\color{white}
	 \int d^4x  \int d^4\theta \, \mathrm{sdet}(E_M{}^A) (\bar GG)^{1/3} } {\color{Orange} \sqrt{g(F)} }\, {\color{TealBlue}f( x)}
\nonumber
\end{align} 
}
is uniquely fixed by \cite{Becker:2016xgv, Becker:2016rku, Becker:2017njd}
\begin{itemize}
	\item Local supersymmetry on 4D, $N=1$, 
	\item tensor hierarchy gauge invariance, 
	\item {\color{Orange} diffeomorphism invariance in 7D}, 
	\item {\color{teal} engineering dimensions, 4D conformal weights, $U(1)_R$ weights, and $GL(7)$ representation} 
	\begin{align}
	{\color {TealBlue} x = {H_i g^{ij}(F) H_j \over (\bar GG)^{2/3}} }
	\nonumber
	\end{align}
\end{itemize}
%up to a function of 
%\begin{align}
%{\color {teal} x = {H_i g^{ij}(F) H_j \over (\bar GG)^{2/3}} }
%\nonumber
%\end{align}
%This is the unique dimensionless, weightless, and chargeless scalar.

Introducing the gravitino superfield $\Psi_i$, we can extend the local supersymmetry from 4D, $N=1$ to all of 11D, $N=1/8$. This fixes $f(x)$ exactly: 
\begin{align}
	\tfrac x4 {\hat f}^4 +  {\hat f}^3 -1 = 0
		~~~\with~~~
	{\hat f} := f -2 x f'	
\nonumber
\end{align}
and the boundary condition $f(0)= 1$ $\Rightarrow$ $f(x) = $ 
%\textrm{ugly but explicit function.}
ugly but explicit function.

%The combination ${\hat f} := f -2 x f'$ satisfies the quartic equation 

\begin{figure}[h]
\center
\scalebox{0.40}{
\includegraphics{./Butter.png}
}
\begin{caption}{
The function $f(x)$ appearing in the D-term integral interpolates starts as $1+\tfrac x{12}$ and slowly asymptotes to $\sqrt{x}/3$.}
\label{F:Fhat}
\end{caption}
\end{figure}

%%%%%%%%%%%%%%%%%%%%%%%%%%%%%%%%%%%%%%%%%%%%%%%%%%%%%%%%%%%%%%%%
\subsection{Components}
A big mystery is how in the world this can produce the correct ``scalar potential'' 
\begin{align}
\label{E:ScalarPotential}
V(x,y)&= 
	-\tfrac12 R(g)
		+\tfrac1{4\cdot 4!} F_{ijkl}^2 
\tag{$\ast \ast$}
\end{align}
$=$ Einstein-Hilbert and Maxwell terms with all polarizations along 7D.\\
In hybrid superspaces, this must happen by (more complicated) auxiliary field mechanism: 
In terms of these fields, the scalar potential is 
\begin{align}
%\label{E:ScalarPotential1}
\hspace{-8mm}V=
	\tfrac14d_X^2 
	+ 2 g_{ij} \bm d^i \bm d^j 		
	-  d_{\mathbf 7 ij} ^2 
		+ \tfrac12 d_{\mathbf {14} ij}^2
	-\tfrac1{18} |f_{\mathbf 1 ijk}+\tfrac{i}2d_XF_{ijk}|^2
		-\tfrac1{24} |f_{\mathbf 7 ijk}|^2
		+\tfrac1{24} |f_{\mathbf {27} ijk}|^2
\nonumber
\end{align}
Define the ``torsion forms'' of $\varphi$ by 
\begin{align}
d\varphi = \tau_0 \ast \varphi + 3\tau_1 \wedge \varphi + \ast \tau_3
~~~\and~~~
d(\ast \varphi) = 4\tau_1 \wedge \ast \varphi + \tau_2 \wedge \varphi
\nonumber
\end{align}
and define $\sigma_0,\cdots, \sigma_3$ similarly for the 3-form $C_{ijk}$.\\
Then component projection results in the equations of motion for the auxiliary fields. Plugging them in gives \cite{Becker:2016edk}
\begin{align}
%\label{E:ScalarPotential}
V(x,y)&= 
	-\tfrac{21}{16} \tau_0^2 
	- 15\tau_1^2 
	+\tfrac18 \tau_2^2 
	+\tfrac1{24} \tau_3^2			
%
	+\tfrac{7}{4} \sigma_0^2 
	+ 9 \sigma_1^2  
	+\tfrac1{24} \sigma_3^2	
=\eqref{E:ScalarPotential}	
\nonumber
\end{align}


%%%%%%%%%%%%%%%%%%%%%%%%%%%%%%%%%%%%%%%%%%%%%%%%%%%%%%%%%%%%%%%%
\subsection{Variational equations}

Variation of the action with respect to prepotentials
{%\color{white}
\begin{align}
\delta S = \int 
\Big\{
	\delta U^a E_{U^a}
	+\delta \Psi_i^\alpha E_{\Psi_i^\alpha}
	+\delta \mathcal V^i E_{\mathcal V^i}
	+\delta X E_{X}
	+\delta \Sigma_i^\alpha E_{\Sigma_i^\alpha}
	+\delta V_{ij} E_{V_{ij} }
	+\delta \Phi_{ijk} E_{\Phi_{ijk} }
	\Big\}
\nonumber	
\end{align}
}
defines the supercurrents
{%\color{white}
\begin{subequations}
\label{E:EoM}
\begin{align}
-9 E_{U^{a}}  &= -18 \dX_a + \varphi^{ijk} \G_{a\, ijk}
\cr
4 E_{\Psi_i^\alpha} &= \l_i^\alpha
\cr
6E_{\mathcal V^i} &= - 6\dY_i
	 +\varphi^{jkl}  \G_{ijkl}
\cr
96 E_X&= -128 \d
 	-i\varphi_{ijk} (\f^{ijk} -\bar \f^{ijk})
	+\psi^{ijkl} \G_{ijkl}
\cr
-16E_{\Sigma^\alpha_i} &= 
	\bar D^2 \l_\alpha^i
\cr
E_{V_{ij} } &= - \psi^{ijkl} \dYY_{kl} 
\cr
-48 E_{\Phi_{ijk} } &=  \f^{ijk} 
	-\tfrac{8i}3\varphi^{ijk} \d 
	+ 18\varphi^{ijk}  \varepsilon^{abcd} \G_{abcd} 
	+12i \varepsilon^{ijklmnp} \G_{lmnp}
%~.
\nonumber
\end{align}
\end{subequations}
}
Notes
\begin{itemize}
	\item These superfields satisfy differential relations expressing the pre-gauge symmetry (redundancy) in the superfield description. For example $E_\Sigma = -\tfrac14 \bar D^2 E_\Psi$ because both fields shift under a chiral spinor parameter superfield: $\delta \Psi_i = \Xi_i+D\Omega + \cdots $ and $\delta \Sigma = -\Xi+... $ compensates. 
	\item Going on-shell corresponds to Imposing equations of motion $E_\ast \to 0$.
	\item On-shell every auxiliary field either vanishes or becomes related to a component of the 4-form field strength $\G = dC$.
\end{itemize}



%%%%%%%%%%%%%%%%%%%%%%%%%%%%%%%%%%%%%%%%%%%%%%%%%%%%%%%%%%%%%%%%
%%%%%%%%%%%%%%%%%%%%%%%%%%%%%%%%%%%%%%%%%%%%%%%%%%%%%%%%%%%%%%%%
\section{Supergeometry}

We can now read off the low-dimension field strength, torsions, and curvatures of the 11D, $N=1/8$ supergeometry.
\begin{itemize}
\item There are no dimension-0 invariants, so 
{%\color{white}
\begin{gather}
\Torsion_{\alpha \beta}{}^{c} = 0
~,~~
\Torsion_{\alpha \dt \beta}{}^{c} = -2i (\sigma^c)_{\alpha \dt \beta}
~,~~
\Torsion_{\hat \alpha \hat \beta}{}^{ k} = 0
%~,~~
\nonumber
\end{gather}
}
where $\hat \alpha = \alpha, \dt \alpha$.
\item At dimension 1/2, there is only the gravitino equation of motion. 
With our minimal choice, it only enters as the dimension-$1/2$ component of the 4-form
{%\color{white}
\begin{gather}
\G_{\alpha b ij} = \varphi_{ijk} (\sigma_b\bar \l^{k})_\alpha  
%~.
\nonumber
\end{gather}
}
This allows for the simplest set of dimension-$\frac12$ torsions 
{%\color{white}
\begin{gather}
%\Torsion_{\hat \alpha \hat \beta}{}^{\hat \gamma} = 0
%~,~~
%\Torsion_{\hat \alpha \hat \beta}{}^{\gamma k} = 0
\Torsion_{\hat \alpha \hat \beta}{}^{\bm \gamma} = 0
~,~~
%\Torsion_{\alpha b}{}^{c} = 0
%~,~~
%\Torsion_{\alpha j}{}^{c} = 0
%~,~~
%\Torsion_{\alpha b}{}^{k} = 0
%~,~~
%\Torsion_{\alpha j}{}^{ k} = 0
\Torsion_{\alpha \bm b}{}^{\bm c} = 0
%~.
\nonumber
\end{gather}
}
for {%\color{white} 
$\bm \gamma = \hat \gamma, \hat \gamma k$ and $\bm c = c, k$}
\item The spin connections appearing in the supersymmetry transformations imply the dimesion-1 conventional constraints
{%\color{white}
\begin{gather}
\Torsion_{\bm{a b}}{}^{\bm c} = 0
%\Torsion_{a b}{}^{c} = 0
%~,~~
%\Torsion_{a j}{}^{c} = 0
%~,~~
%\Torsion_{i j}{}^{c} = 0
%~,~~
%\Torsion_{a b}{}^{k} = 0
%~,~~
%\Torsion_{a j}{}^{ k} = 0
%~,~~
%\Torsion_{i j}{}^{ k} = 0
\end{gather}
}
\item With this, we can proceed to solve the superspace Bianchi identities.
{%\color{white}
\begin{align}
d \G + \iota_\Torsion \G = 0
~~~\and~~~
d\Torsion +\iota_\Torsion \Torsion = \Curvature 
\nonumber
\end{align}
}
for the remaining field strength, torsion, and curvature components (Dragon's theorem). 
\end{itemize}

For example, the dimension-1 torsions with $N=1$ spinor index are
\begin{gather}
\Torsion_{\alpha b}{}^{\gamma} = 
	\tfrac i2 (\delta_b^c - \sigma_b \bar \sigma^c)_{\alpha}{}^\gamma  \dX_c
~,~~
\Torsion_{\alpha b}{}^{\dt \gamma} = 
	- i (\sigma_b)_{\alpha}{}^{\dt \gamma} {\bm R}
~,~~
%	\cr
%T_{\alpha b}{}^{\gamma k} =
%	-\tfrac i{12} (\delta_b^c - \sigma_b \bar \sigma^c)_{\alpha}{}^\gamma \bar \Ga_c^k
%	-\tfrac i{16} \delta_\alpha^\gamma (\sigma_b)^{\un b} \varphi^{ijk} R_{\beta \dt \beta i j}
%~,~~
%\cr
%T_{\alpha b}{}^{\dt \gamma k} =
%	i(\sigma_b)_\alpha{}^{\dt \gamma} \d^k 
%	-\tfrac1{24} (\sigma_b)^{\un c} \t_{\alpha \gamma}{}^k
%	-\tfrac1{16} (\sigma_b)_{\un a} \bar \t^{\dt \alpha \dt \gamma k}
%~,~~
	\cr	
\Torsion_{\alpha j}{}^{\gamma} = 
	\tfrac12 \delta_\alpha^\gamma \d_j 
	+\tfrac1{48} (\sigma^{ab})_\alpha{}^\gamma \t_{ab \,j }
~,~~
\Torsion_{\alpha j}{}^{\dt \gamma} = -\tfrac i{12} (\sigma_a)_\alpha{}^{\dt \gamma} \bar \Ga_j^a
~,~~
%	\cr
%T_{\alpha j}{}^{\gamma k} \sim \tfrac12 \delta_\alpha^\gamma Z_j{}^k
%	-\tfrac i2 (\sigma^{ab})_\alpha{}^\gamma \pi_{\bm {14}}\F_{ab\, j}{}^k 
%~,~~
%T_{\alpha j}{}^{\dt \gamma k} \sim (\sigma^a)_{\alpha}{}^{\dt \gamma} \varphi^{klm} \I_{a jlm}
%~,~~	
\end{gather}
whereas the dimension-1 torsions with $\bm 7$-valued spinor index are
\begin{gather}
%T_{\alpha b}{}^{\gamma} = 
%	\tfrac i2 (\delta_b^c - \sigma_b \bar \sigma^c)_{\alpha}{}^\gamma  \dX_c
%~,~~
%T_{\alpha b}{}^{\dt \gamma} = 
%	- i (\sigma_b)_{\alpha}{}^{\dt \gamma} {\bm R}
%~,~~
%	\cr
\Torsion_{\alpha b}{}^{\gamma k} =
%	-\tfrac i{12} (\delta_b^c - \sigma_b \bar \sigma^c)_{\alpha}{}^\gamma \bar \Ga_c^k
%	-\tfrac i{16} \delta_\alpha^\gamma (\sigma_b)^{\un b} \varphi^{ijk} R_{\beta \dt \beta i j}
-\tfrac i{12} (\delta_b^c - \sigma_b\bar \sigma^c)_\alpha{}^\gamma \bar \Ga_c^k 
	+ \tfrac 18 \delta_\alpha{}^\gamma \check \Ga_b^k 
	-\tfrac i{32} \delta_\alpha^\gamma 
		[ \bar D \bar \sigma_b \l^k + D\sigma_b \bar \l{}^k]
~,~~
\cr
\Torsion_{\alpha b}{}^{\dt \gamma k} =
	i(\sigma_b)_\alpha{}^{\dt \gamma} \d^k 
	-\tfrac1{24} (\sigma_b)^{\un c} \t_{\alpha \gamma}{}^k
	-\tfrac1{16} (\sigma_b)_{\un a} \bar \t^{\dt \alpha \dt \gamma k}
~,~~
	\cr	
%T_{\alpha j}{}^{\gamma} = 
%	\tfrac12 \delta_\alpha^\gamma \d_j 
%	+\tfrac1{24} \t_{j \alpha}{}^\gamma
%~,~~
%T_{\alpha j}{}^{\dt \gamma} = -\tfrac i{12} (\sigma_a)_\alpha{}^{\dt \gamma} \bar \Ga_j^a
%~,~~
%	\cr
\Torsion_{\alpha j}{}^{\gamma k} \sim \tfrac12 \delta_\alpha^\gamma \bar \Oh_j{}^k
	-\tfrac i2 (\sigma^{ab})_\alpha{}^\gamma \pi_{\bm {14}}\G_{ab\, j}{}^k 
~,~~
\cr
\Torsion_{\alpha j}{}^{\dt \gamma k} = 
	\tfrac i4 (\sigma^a)_\alpha{}^{\dt \gamma} [
		h_j{}^k(\I_{a}) 
	-\tfrac 1{6} \varphi_j{}^{kl} \bar \Ga_l^a
	]
	+\tfrac i{24} \varphi_j{}^{kl} D_\alpha \bar \l^{\dt \gamma}_l
~.~~	
\end{gather}
%The dimension-$\frac32$ curls
%\begin{align}
%T_{a b}{}^{\gamma} = \Weyl_{a b}{}^{\gamma} 
%	+ (\bar D \tilde \sigma_{[a} )^{\gamma} \dX_{b]}
%~,~~
%T_{a b}{}^{\gamma k} = (d\physical \psi)_{a b}{}^{\gamma k} 
%~,~~
%	\cr
%T_{a j}{}^{\gamma} = (d\physical \psi)_{a j}{}^{\gamma} + D\Ga
%~,~~
%T_{a j}{}^{\gamma k} =  (d\physical \psi)_{a j}{}^{\gamma k}
%~,~~
%	\cr
%T_{i j}{}^{\gamma} = (d\physical \psi)_{i j}{}^{\gamma} + D\F
%~,~~
%T_{i j}{}^{\gamma k} = (d\physical \psi)_{i j}{}^{\gamma k}
%\end{align}
and so on.
These results agree with our supersymmetry transformations 
\begin{align}
\delta E_M{}^{\bm C} = - \mathcal D_M \xi^{\bm C} - \xi^{\bm D} T_{\bm D M}{}^{\bm C}
\end{align}


%%%%%%%%%%%%%%%%%%%%%%%%%%%%%%%%%%%%%%%%%%%%%%%%%%%%%%%%%%%%%%%%
%%%%%%%%%%%%%%%%%%%%%%%%%%%%%%%%%%%%%%%%%%%%%%%%%%%%%%%%%%%%%%%%
\section{Toward the $R^2$ Invariant}

{%\color{white}
\begin{gather}
\Curvature_{\bm{ab}}{}^{\bm {cd}} 
\cr
\TorsionX_{\bm {ab}}{}^\gamma,
\TorsionY_{\bm {ab}}{}^{\gamma k},
%\CurlXXX_{\alpha \beta}{}^{ \gamma} ,
%\CurlXYX_{aj}{}^{\gamma} , 
%\CurlYYX_{ij}{}^{\gamma} , 
%\CurlXXY_{\alpha \beta}{}^{\gamma k} , 
%\CurlXYY_{aj}{}^{\gamma k} , 
%\CurlYYY_{ij}{}^{\gamma k} , 
\r_i^\alpha 
\cr
\dX^a, 
\d, 
\dY^i, 
\dYY_{ij}, 
\f_{ijk},
\y_i^a,
\t_k^{ab},
\G_{\bm {abcd}}
\cr
\l_i^\alpha
\nonumber
\end{gather}
}

{%\color{white}
\begin{align}
\mathbb G &= \Weyl^{\alpha \beta \gamma}\Weyl_{\alpha \beta \gamma}
	-\tfrac14 \bar D^2 (\dX_a^2)
	+\cdots	
\cr
%\label{E:mathbbH}
\mathbb H_i &= 
%-\tfrac12 D^\gamma F_i^{\prime \alpha \beta} W_{\alpha \beta \gamma}
	4i \Weyl_i^{\alpha \beta \gamma} \Weyl_{\alpha \beta \gamma}
	-\tfrac 14 D^{\alpha} (\bar F_i^{\prime \dt\alpha \dt\beta}\bar D_{\dt\beta} G_{\alpha \dt \alpha} -\tfrac18 D^2 \lambda_{\alpha i} R)
\cr	&+\tfrac 1{24} \left( D^2 \bar D_{\dt \alpha}\lambda_{\alpha i}+ 2 D^\beta \bar D_{\dt \alpha} D_{(\alpha}\lambda_{\beta)i} \right) \, G^{\alpha \dt \alpha} 
	+\hc+\cdots	
\cr
\mathbb G &= \Weyl^{\alpha \beta \gamma} \Weyl_{\alpha \beta \gamma} 
\cr
\mathbb H_i &= D^\gamma \KK_i^{\alpha \beta} \Weyl_{\alpha \beta \gamma}
	-\tfrac 12 \KK_i^{\alpha \beta}\bar D^{\dt \alpha}  D_\beta \dX_{\alpha \dt \alpha}
	-\tfrac 12 \bar D^{\dt \alpha} \KK_i^{\alpha \beta}D_\beta \dX_{\alpha \dt \alpha} 
\cr	&	-\tfrac 1{4} \left( D^\beta \bar D_{\dt \alpha} \KK_{\alpha \beta i} 
	+\tfrac 14 D_\alpha \bar D^2 \bar \l_{\dt \alpha i} \right) \dX^{\alpha \dt \alpha}
	+\tfrac 1{32} D^2 \l_i^\alpha D_\alpha \Scalar
	+\cdots
	+ \mathrm{h.c.}
\cr
\mathbb W_{\alpha ij} &= \tfrac i8 R_{ij}{}^{\beta \gamma} \Weyl_{\alpha \beta \gamma}
	+2i \partial_{[i}\mathbb Z_{\alpha j]} 
	+ i \bar \Theta^{\dt \alpha \dt \beta \dt \gamma}_{[i} D_\alpha \bar \Theta_{j]\, \dt \alpha \dt \beta \dt \gamma}
	-2i \bar R_{ij}{}^{\dt \alpha \dt \beta} \bar D_{\dt \beta} G_{\alpha \dt \alpha}
	+\cdots	
\cr
\mathbb F_{ijk} &= b R_{[ij}{}^{\alpha \beta} F'_{k]\alpha \beta}
	+\cdots	
\cr
%\label{E:mathbbE}
\mathbb E_{ijkl} &= a R_{[ij}{}^{\alpha \beta} R_{kl] \alpha \beta} +\cdots 	
\nonumber	
\end{align}
}
In this form, $\mathbb H_y'$ is manifestly real again. 

Another form in which it manifestly satisfies the descent relation is
{\color{white}
\begin{align}
%\label{E:mathbbG}
\mathbb G &= \TorsionXXX^{\alpha \beta \gamma} \TorsionXXX_{\alpha \beta \gamma}
	-\tfrac14 \bar D^2 (\dX_a^2)
	+\cdots	
\cr
%\label{E:mathbbH}
\mathbb H_i &= 
%-\tfrac12 D^\gamma F_i^{\prime \alpha \beta} W_{\alpha \beta \gamma}
%	+\tfrac 14 F_i^{\prime \alpha \beta}\bar D^{\dt \alpha}  D_\beta G_{\un a}
%	+\tfrac 14 \bar D^{\dt \alpha} F_i^{\prime \alpha \beta}D_\beta G_{\un a} 
%\cr	&+\tfrac 1{24} \left( D^\alpha \bar D^{\dt \alpha} + 2 \bar D^{\dt \alpha} D^\alpha \right) 
%D^\beta \lambda_{\beta \,i} \, G_{\un a} 
%	+\tfrac 1{32} D^2 \lambda_i^\alpha D_\alpha R
%	+ ~\hc
%%	- \partial_i(G_a^2)
%	+\cdots
-2i \TorsionXXY_i^{\alpha \beta \gamma} \TorsionXXX_{\alpha \beta \gamma}
	- \partial_i(\dX_a^2)
	+\bar D_{\dt \alpha} \bar {\mathbb Z}_i^{\dt \alpha}
	+\cdots	
\cr
%\label{E:mathbbW}
\mathbb W_{\alpha ij} &= \tfrac i8 \Curvature_{ij}{}^{\beta \gamma} \TorsionXXX_{\alpha \beta \gamma}
	+2i \partial_{[i}\mathbb Z_{\alpha j]} 
	+ i \bar \TorsionXXY^{\dt \alpha \dt \beta \dt \gamma}_{[i} D_\alpha \bar \TorsionXXY_{j]\, \dt \alpha \dt \beta \dt \gamma}
%	-2i \bar \Curvature_{ij}{}^{\dt \alpha \dt \beta} \bar D_{\dt \beta} \dX_{\alpha \dt \alpha}
	+\cdots	
\cr
\mathbb F_{ijk} &= c_3 \Curvature_{[ij}{}^{\alpha \beta} \KK_{k]\alpha \beta}
	+\cdots	
\cr
%\label{E:mathbbE}
\mathbb E_{ijkl} &= c_4 \Curvature_{[ij}{}^{\alpha \beta} \Curvature_{kl] \alpha \beta} +\cdots 	
\nonumber 	
\end{align}
}
with
{%\color{white}
\begin{align}
%\label{E:mathbbZ}
\bar {\mathbb Z}_i^{\dt \alpha}  &:= 
-\tfrac12 \bar \KK_{\dt \beta \dt \gamma \,i} \bar \TorsionXXX^{\dt \alpha \dt\beta \dt\gamma}
-\tfrac 14 \KK_{\alpha \beta\,i}D^\beta \dX^{\alpha \dt \alpha} 
+\tfrac 1{4} D^{\beta} \bar \KK_i^{\dt \alpha \dt \beta}  \dX_{\beta \dt \beta}
	+\tfrac i{32} D^2 \l_{\alpha i} \dX^{\alpha \dt \alpha} 
	+\tfrac i{32} \bar D^2 \bar \l_i^{\dt \alpha} \bar \Scalar
	+\cdots
\cr
\KK_{ab}{}^k &:= -\tfrac{4i}3 \varphi^{ijk} \G_{ab\, jk} + \tfrac12 D\sigma_{ab} \l^k
\nonumber	
\end{align}
}

{%\color{white}
\begin{align}
%\bar D_{\dt \alpha} \mathbb G = 0
%~,~
\bar D^2 \mathbb H =  4\partial \mathbb G 
~,~
D\mathbb W - \bar D\bar {\mathbb W}= 2i\partial \mathbb H 
~,~
\bar D^2 D\mathbb F  = 4\partial \mathbb W 
~,~
\mathbb E - \bar {\mathbb E} &=  2i \partial \mathbb F
~,~
\partial \mathbb E = 0
\nonumber
\end{align}
}

%%%%%%%%%%%%%%%%%%%%%%%%%%%%%%%%%%%%%%%%%%%%%%%%%%%%%%%%%%%%%%%%
%%%%%%%%%%%%%%%%%%%%%%%%%%%%%%%%%%%%%%%%%%%%%%%%%%%%%%%%%%%%%%%%
\section{Beyond...}


%%%%
%%%% Tabel 3
%%%%
%\begin{table}[t]
{\color{white}
\footnotesize
\begin{align*}
~\hspace{-13mm}
{\renewcommand{\arraystretch}{1.5} %adds some padding
%\begin{array}{|c|ccccccc|}
\begin{array}{|c|c|c|c|c|c|c|c|c|}
\hline
& U^a & \Psi_i^\alpha & \mathcal V^i  & X & \Sigma_{\alpha i} & V_{ij} &  \Phi_{ijk} &\textrm{sdim}_{\mathbf R} \\
\textrm{prepotential}& \textrm{real} & \textrm{spinor} & \textrm{real} &  \textrm{real} &  \textrm{chiral}  &  \textrm{real} &  \textrm{chiral}&\\
\hline
\physical{\textrm{physical }}
	&\gravitonXX_{ab} , \grinoXX_a{}^\beta 
	& \grinoXY_a{}^{\beta j} & \gravitonXY_a^i & \CXXX_{abc} 
	& \CXXY_{ab\,i} & \CXYY_{a\, ij}, \spinorYY_{\alpha\, ij} 
	& \CYYY_{ijk}, \gravitonYY_{ij}, \spinorYYY_{\alpha\, ijk}&\\
%	& 2|2 & \\
\auxiliary{\textrm{auxiliary }} & \dX^a 
	&\l^i_\alpha, \y_i^a, \t_i{}^{\alpha \beta} ,\r_i^\alpha 
	& \dY^i & \d &  & \dYY_{ij} & \f_{ijk} &\\
\compensating{\textrm{compensating}}
	&&&\surfeit^i_\alpha &\compX,\surfeit'_\alpha 
	&\compY_i,\surfeit''_{\alpha i}&&&
	\\
\hline	
\textrm{on-shell} & 2|2 & 0|14 & 14|0 & 0|0 & 7|0 & 42|42 & 63|70 & 128|128 \\
%\textrm{off-shell} & 10|12 & 98|140 & 28|0 & 2|0 & 21|0 & 84|84 & 133|140 &   376|376\\
\textrm{off-shell} & (\physical{6}+\auxiliary{4})|\physical{12} & \auxiliary{98}|(\physical{84}+\auxiliary{56}) & (\physical{21}+\auxiliary{7})|0 & (\physical{1}+\auxiliary{1})|0 & \physical{21}|0 & (\physical{63}+\auxiliary{21})|\physical{84} & (\physical{63}+\auxiliary{70})|\physical{140} &   376|376\\
\hline
\end{array}
}
\end{align*}
%\caption{\scriptsize The unconstrained ``prepotential'' superfields of the 11D, $N=1/8$ supergravity multiplet and their physical, auxiliary, and compensating component content. The components of the gravitino of the type $\psi_i{}^\beta$ and $\psi_i{}^{\beta j}$ are encoded in $\chi_{\alpha ij}$ and $\chi_{\alpha ijk}$ \eqref{E:spinor}. The real superdimension contributed by each superfield is computed modulo gauge transformation. The superdimension of the auxiliary field space adds up to $201|56$. 
%%The total off-shell dimension of this supermultiplet is 
%}
%\label{T:Components}
} %end scriptsize
%\end{table}
%%%%
%%%% Tabel 3 
%%%%


{\color{white}
\begin{align*}
\delta_\epsilon \grinoXX_a^\beta &\simeq
	-i  (\epsilon\delta_a^b - \epsilon \sigma_a \bar \sigma^b)^\beta \dX_b
	-\tfrac i2 (\bar \epsilon \bar \sigma_a )^{\beta} \Scalar
\\
%%%%%
\delta_\epsilon \grinoYX_i^\beta &\simeq
	\tfrac i{12} (\epsilon \sigma^{ab})^\beta \t_{ab\, i}
	+\tfrac i2 \epsilon^\beta \dY_i  
	+ \tfrac i{12} (\bar \epsilon \bar \sigma_a )^\beta \Ga_i^{a}
\cr
%%%%%
\delta_\epsilon \grinoXY_{a}{}^{\beta j} 
	&\simeq 
	\tfrac{5i}{12} (\epsilon\delta_a^b - \epsilon \sigma_a \bar \sigma^b)^\beta\,  \bar \Ga_b^j
	-i (\bar \epsilon\bar \sigma_b )^\beta
		\big[ 
			\delta_a^b \dY^j 
			-\tfrac i2 \varepsilon_a{}^{bcd} \t^j_{cd} 
		\big]
%\cr&
+\tfrac18 \epsilon^\beta \big[
		2i \Ga^j_a
		-\bar D\bar \sigma_a \l^j + D\sigma_a \bar \l^j 
	\big]
\cr
%%%%%
\delta_\epsilon \grinoYY_{i}{}^{\beta j} &\simeq
	-i(\epsilon \sigma^{ab})^\beta \pi_{\bm {14}} \FSXXYY_{ab\,i}{}^j
	-i (\bar \epsilon \bar \sigma^a)^\beta \left[
		h_i^j(\FSXYYY_{a\, ijk})
		+\tfrac i{12}\varphi_i{}^{jk} ( 
			2i  \Ga_{a k}
			-\bar D \bar \sigma_a \l_k
			)
		\right]
	+2\epsilon^\beta \Ztorsion_i{}^j
\nonumber	
\end{align*}
}




{\color{white}
\begin{align*}
%\hyperlabel{E:R}
%\R
\R &:= \tfrac13 \d + \tfrac i{6\cdot 4!} \varepsilon^{abcd} \FSXXXX_{abcd}
\\
%\hyperlabel{E:Ga}
\Ga^{a}_i 
	&:=  \tfrac 1{2i} \big[ \y_i^{a} - \bar \y_i^{a} \big]
	+\tfrac16 \psi_{ijkl} \FSXYYY^{a jkl} 
	+\tfrac i{3!} \varepsilon^{abcd} \FSXXXY_{bcd\, i}
\\
%\hyperlabel{E:Ztorsion}
\Ztorsion_i{}^j 
	&:= 
	- i h_i{}^j ( \f)
	+\tfrac1{36}\varphi_i{}^{jk}\psi_k{}^{lmn} \f_{lmn}
	+  \pi_{\bm{14}} \dYY_i{}^j	
\end{align*}
}




{\color{white}
\begin{align}
\dX_a &=\tfrac1{18} \varphi^{ijk} \FSXYYY_{a\, ijk} 
	+ \tfrac12 \EoMX{}_a
\\
\l_\alpha^i &= 4 \EoMgrino{}_\alpha^i
\\
\dY_i&= \tfrac16\varphi^{jkl} \FSYYYY_{ijkl}
	 - \EoMY{}_i
\\
\d&= 
	-\tfrac 1{96}\psi^{ijkl}\FSYYYY_{ijkl}
	+\EoM
	-\tfrac i2 \varphi_{ijk} ( \EoMYYY - \overline {\EoMYYY} )^{ijk}
\\
\dYY_{ij}  &= \tfrac1{8} \left(
	2 \delta_{[i}^k \delta_{j]}^l + \psi_{ij}{}^{kl}
	\right)
	\EoMYY{}_{kl} 
\\
%\f^{ijk}  &= 
%- \tfrac 1{3\cdot 4!}\left[
%	4 \varphi^{ijk} \varepsilon^{abcd} \FSXXXX_{abcd} 
%	+i \varphi^{ijk} \psi^{lmnp} \FSYYYY_{lmnp} 
%%	+6i \epsilon^{ijklmnp} \FSYYYY_{lmnp}
%	\right]-48\EoMYYY^{ijk}
%\cr
%	&\hspace{45mm}
%	+\tfrac{4}3 \left[
%		\varphi^{ijk} \varphi_{lmn} ( \EoMYYY - \overline {\EoMYYY})^{lmn}
%		+2i\EoM 
%	\right]
G^{ijk\, lmn} \f_{lmn}  &= 
\tfrac 1{36}\left[
	\varphi^{ijk}( 2i  \varepsilon^{abcd} \FSXXXX_{abcd} 
	+  \psi^{lmnp} \FSYYYY_{lmnp} )
%	2i \varphi^{ijk} \varepsilon^{abcd} \FSXXXX_{abcd} 
%	+ \varphi^{ijk} \psi^{lmnp} \FSYYYY_{lmnp} 
	+3 \varepsilon^{ijklmnp} \FSYYYY_{lmnp} 
	\right]
\cr
	&\hspace{25mm}
	+48i\overline\EoMYYY^{ijk}
	+\tfrac{4i}3 \varphi^{ijk} \varphi_{lmn} ( \EoMYYY - \overline {\EoMYYY})^{lmn}
	-\tfrac 83\EoM 
\end{align}
}

{%\color{white}
\begin{align}
%\y_{\un a i} &= - 4\bar D_{\dt \alpha} \EoMgrino{}_\alpha^i
\y_{i}^a&= 2 \bar D \bar \sigma^a \EoMgrino{}_i
\cr
\t_{\alpha \beta i} &=
		-\tfrac12 \varphi_i{}^{jk} \FSXXYY_{\alpha \beta jk} 
		-i D_{(\alpha}\EoMgrino{}_{\beta) i}
\cr
\r_i^\alpha 
%	&= -\tfrac 1{12} \psi_i{}^{jkl} \bar D_{\dt \alpha} \FSXYYY_{jkl}^{\alpha \dt \alpha}
%	+\tfrac i2 \bar D_{\dt \alpha} \tilde \FSXXXY_i^{\un a}
%	+2i \psi_i{}^{jkl} \partial_j \spinorYY_{kl}^\alpha 
%	- i \bar D_{\dt \alpha} D^\alpha \overline \EoMgrino^{\dt \alpha}_{ i} 
%\cr
	&=
	-\tfrac 1{12} (\bar D\bar \sigma^a)^{\alpha} 
		\left[ 
			\psi_i{}^{jkl} \FSXYYY_{a\, jkl} 
			-i \varepsilon_a{}^{bcd}  \FSXXXY_{bcd\, i}
		\right]
	+2i \psi_i{}^{jkl} \partial_j \spinorYY_{kl}^\alpha 
	- i \bar D_{\dt \alpha} D^\alpha \overline \EoMgrino^{\dt \alpha}_{ i} 
\end{align}
}


{%\color{white}
\begin{displaymath}
\xymatrix{
	&
	\Scalar \ar[dl]_{\bar D} \ar[dr]^{D} 
	&& 
	\dX_a \ar[dl]_{\bar D\cdot } \ar[dr]^{\bar D\wedge}  
\\
	0
	&& 
	\Torsion_\alpha \ar[dl]_{D} \ar[dr]^{\bar D} 
	&& 
	\TorsionXXX_{\dt \beta \dt \alpha}{}^\gamma \ar[dl]_{\bar D} \ar[dr]^{D}
	&&
	\TorsionXXX_{\alpha \beta}{}^\gamma\ar[dl]_{D} \ar[dr]^{\bar D}
\\
	&
	\Curvature_{ab}{}^{ab}
	& & 
	\partial \Curvature_{ab}{}^{ab}
	&&
	\Curvature_{\dt \alpha\dt \beta}{}^{\gamma \delta} \oplus \partial \dX^a \oplus
	\Curvature_{\alpha \beta}{}^{\gamma \delta} 
	&&
	0
}
\end{displaymath}
}

{%\color{white}
{\renewcommand{\arraystretch}{1.7} %adds some padding
\begin{align*}
\begin{array}{llll}
T_{\alpha b}{}^{\gamma} \sim \dX_b
&
T_{\alpha b}{}^{\dot \gamma} \sim  \R
&
T_{\alpha j}{}^{\gamma} \sim  \dY_j  ~,~ \t_{ab \,j}
&
T_{\alpha j}{}^{\dot \gamma} \sim \bar \Ga_j^a
\\
T_{\alpha b}{}^{\gamma k} \sim \Ga_b^k~,~ \bar D \bar \sigma_b \l^k
&
T_{\alpha b}{}^{\dot \gamma k} \sim \dY^k~,~ \t_{ab}{}^k
&
T_{\alpha j}{}^{\gamma k} \sim \Ztorsion_j{}^k~,~ \FSXXYY_{ab\, [jk]_{\bm {14}}}
&
T_{\alpha j}{}^{\dot \gamma k} \sim \FSXYYY_{a\,ijk}~,~\Ga_{a i}~,~D_\alpha \bar \l^{\dt \gamma}_i
\end{array}
\end{align*}
}
}




\center
\scalebox{0.40}{
\includegraphics{./5DButter}
\includegraphics{./11DButter}
}








{\color{white}
{\renewcommand{\arraystretch}{1.7} %adds some padding
\begin{align*}
\begin{array}{|c|ccccccccc|}
\hline
&\dX^a &\R &\dY^i  & \t_{ab \,k} & \Ga_i^a & \bar D \bar \sigma_a \l^i & \Ztorsion_i{}^j& \FSXXYY_{ab\, [jk]_{\bm {14}}}& \FSXYYY_{a\,[ijk]_{\bm{28}}}
\\
\hline
T_{\alpha b}{}^{\gamma} 
&\ast&&&&&&&&
\\
T_{\alpha b}{}^{\dot \gamma} 
&&\ast&&&&&&&
\\
T_{\alpha j}{}^{\gamma} 
&&&\ast&\ast&&&&&
\\
T_{\alpha j}{}^{\dot \gamma} 
&&&&&\ast&&&&
\\
\hline
T_{\alpha b}{}^{\gamma k}  
&&&&&\ast&\ast&&&
\\
T_{\alpha b}{}^{\dot \gamma k} 
&&&\ast&\ast&&&&&
\\
T_{\alpha j}{}^{\gamma k} 
&&&&&&&\ast&\ast&
\\
T_{\alpha j}{}^{\dot \gamma k} 
&&&&&\ast&\ast&&&\ast
\\
\hline
\end{array}
\end{align*}
}
}

{\color{white}
%%%%%%%%%%
%%%%%%%%%%
\begin{figure}[htb]
\begin{center}
\begin{tikzpicture}[auto, > = {stealth}, decoration = {
    markings,
    mark = at position 0.75 with {\arrow{>}}}]
	
	\draw[fill=black!10!white, draw=black!50!white] (-2,1) -- (2,1) -- (0,-1) -- (-4,-1) -- cycle;
	
	\draw [postaction = {decorate}] (-1, 0) -- (-1, 2);
	\draw [dashed, draw=black!50!white] (-1, 0) -- (-1, -1);
	\draw (-1, -1) -- (-1, -2);
	\draw [draw=black!50!white] (-2.6, -0.35) -- (-2.6, 1.15);
	
	\draw (-3, 0.5) node[right, rotate=90] {$B_{a6}$};
	\draw (0, 1.55) node[right] {$\partial_6 B \,\neq\, 0$};
	\draw (-0.25, 2.2) node[left] {$x^6$};
	\draw (-0.65, 0.5) node[right] {{\color{black}$B_{ab}$}};
%	\draw [->] (0.25, 0) arc (200:270:12pt) arc (90:40:16pt);
%	\draw (0.65, -0.9) node[right] {$\partial_{[\hat a} F_{\hat b \hat c]} \,=\, 0$};
\end{tikzpicture}
\caption{Properties of the Perry-Schwarz 2-form $F$.
\label{fig:2form}}
\end{center}
\end{figure}
%%%%%%%%%%
%%%%%%%%%%
}































%%%%%%%%%%%%%%%%%%%%%%%%%%%%%%%%%%%%%%%%%%%%%%%%%%%%%%%%%%%%%%%
%%%%%%%%%%%%%%%%%%%%%%%%%%%%%%%%%%%%%%%%%%%%%%%%%%%%%%%%%%%%%%%
%%%%%%%%%%%%%%%%%%%%%%%%%%%%%%%%%%%%%%%%%%%%%%%%%%%%%%%%%%%%%%%
%%%%%%%%%%%%%%%%%%%%%%%%%%%%%%%%%%%%%%%%%%%%%%%%%%%%%%%%%%%%%%%
{
%\small
\footnotesize
%\scriptsize
%
\bibliography{/Users/wdlinch3/Dropbox/Rashoumon/LaTeX/BibTex/BibTex}
\bibliographystyle{unsrt}
%\bibliographystyle{amsplain}
%\bibliographystyle{amsalpha}
}

\end{document}  


